\documentclass[a4paper,10pt,twocolumn]{article}

%Layout
% \setlength\columnsep{.8cm}
% \setlength\topmargin{-1.71cm}
% \setlength\textheight{23.4cm}
% \setlength\textwidth{15.8cm}

%Pacotes
\usepackage{amsthm,amssymb,amsmath}
\usepackage[english]{babel}
\usepackage[utf8]{inputenc}
\usepackage{amsfonts}
\usepackage{graphicx}
\usepackage{float}
\usepackage{ragged2e}
\usepackage{fancyhdr}
\usepackage{titlesec}
\usepackage{natbib}
\usepackage[font=small]{caption}
\usepackage[
layout = a4paper,
includehead,
paperheight = 29.7cm,
headheight = 3.41cm,
paperwidth = 21cm,
textwidth = 15.8cm,
textheight = 23.4cm,
columnsep = 0.8cm,
top = 1.25cm,
bottom = 4.2cm,
left = 2.6cm,
right = 2.6cm,
headsep = 0.3cm,
]{geometry}

\setlength{\bibsep}{1pt}

%Fonte Arial
\renewcommand{\rmdefault}{phv}
\renewcommand{\sfdefault}{phv}

\titleformat
{\section} % command
[display] % shape
{\normalfont\fontsize{13}{13}\bfseries} % format
{} % label
{0cm} % sep
{
    \centering
} % before-code
[
] % after-code
\titlespacing*{\section}
{0pt}
{1.5pt}
{1.5pt}
[0pt]

\begin{document}
\pagestyle{fancy}
\fancyhead{}
\fancyhead[L]{\includegraphics[width=3.41cm]{logo_siicusp.jpg}}
\renewcommand{\headrulewidth}{0pt}
\pagenumbering{gobble}

\twocolumn[{
    % Bloco centralizado (ocupa as duas colunas)
    \begin{center}
        {\fontsize{13}{13}\selectfont \textbf{Estimating the duration of random
        walks}}\\[1em]
        {\fontsize{13}{13}\selectfont \textbf{Eduardo Yukio Garrafa Ishihara}}\\[0.3em]
        {\fontsize{13}{13}\selectfont \textbf{Gustavo Silva Garone}}\\[0.9em]
        {\fontsize{13}{13}\selectfont \textbf{Prof. Elisabeti Kira, Ph.D}}\\[0.8em]
        {\fontsize{13}{13}\selectfont Institute of Mathematics and Statistics,
        University of São Paulo}\\[0.3em]
        {\fontsize{10}{10}\selectfont eduardoyukio.ishihara@usp.br}
    \end{center}
    \vspace{1.5em} % espaço antes de começar o corpo do texto
}]

\section{Objectives}

The Gambler’s Ruin is a random walk, one of the most classical models of a
Markovian stochastic process. In this model, a gambler wagers money on each
round, either gaining or losing varying amounts according to a probabilistic
rule. Such models have many applications in statistical modeling and
computational algorithms \cite{ross_markov_2019}. Their behavior has been
thoroughly studied in the simplest form — the gambler wins $\$1.00$
with probability $p$ and loses $\$1.00$ with probability $1-p$ —
particularly concerning the probability of ruin.
\cite{ross_introduction_2019}. This project's aim is to investigate the
expected game duration in more complex random walks by proposing estimators
and evaluating their performance through computational approaches.
We hope through this approach to shine light on interesting and computationally
cheap ways to estimate random walks.

\section{Materials and methods}

We analyzed the expected value and variance of the duration
of the classical Gambler’s Ruin problem \cite{andel_variance_2012} using
finite difference equations systems. We developed a heuristic
estimator based on an analogy with average velocity in Physics. This estimator yielded
accurate approximations for most random walks, both in the simple case and in
models where the gambler may win or lose amounts different from
$\$1.00$ per round. In the fair game case, that is, when the expected gain
per round is zero, the estimator deviated significantly from the theoretical and
simulated values.

For random walks with more general rules, where analytical solutions are not
available, we evaluated this estimator using Monte Carlo simulations, based on
methods proposed by Ritter \cite{ritter_determining_2011}, using the Python
programming language. As we considered broader rules that greatly
extended the simulation times and required more computational power, we
migrated to the Julia language, as recommended by Godoy
\cite{godoy_evaluating_2023}. We further refined the estimator to account for
ties within a round.

Having identified the limitations of this estimator in one dimension, we extended
our study to random walks in $\mathbb{Z}^d$. By solving systems of equations
computationally with symbolic libraries, we obtained analytical results for unit-step
walks with different probabilistic rules. New estimators were analyzed
for the simple and more complex random walks. Parallel programming enabled quick
simulation of random walks in high dimensions.

\section{Results}

The proposed estimator showed efficacy in predicting the duration of
one-dimensional walks, as depicted by Figure 1, and adapts successfully
when the transition probability depends on the gambler's current fortune.

\begin{figure}[H]
    \centering
    \includegraphics[width=7.5cm]{"graficoest-en.pdf"}
    \centering
    \caption{Comparing the estimator with $100\, 000$ simulations per starting
    position}
\end{figure}

In higher dimensions, the expected duration of the random walk must be obtained
numerically due to the system of finite different equation's exponential growth,
rendering a closed form solution intractable. A surprising result was observed
in this environment: once the dimensionality exceeds roughly $d=100$, the
average time to game completion increases almost linearly with each additional
dimension as shown in Figure 2.

\begin{figure}[h]
    \centering
    \includegraphics[width=7.5cm]{"boxplots-en.pdf"}
    \centering
    \caption{Game duration in $1\, 000$ simulations by dimension.}
\end{figure}

In the particular model where the gambler cannot lose his money, we've obtained
analytic expressions for the random and non-random variants. The complexity of
the equations that describe the game's duration grows proportionally to the
state space's size. In all cases, the estimated values show great proximity to
simulated ones, although greater discrepancies show when the game is close to
fair - that is, when the gambler's expected earnings per round are close to $0$
- or when the starting position is close to a boundary value.

\section{Conclusions}

The Gambler's Ruin is a classical math problem proposed almost four centuries
ago. Through intuition, undergraduates can formulate fresh questions and
hypothesis, explore them computationally, and finally developg rigorous
mathematics. It is our hope that this work encourages other undergraduates to
explore and reexplore classical math problems in search of solutions to modern
day challenges.

The authors declare no conflict of interests. Eduardo Yukio conceived the
theoretical framework behind the estimators and derived the paper's analytical
results. Gustavo Garone provided computational support and implemented the
simulations used for improving the estimators.

\section{Acknowledgements}

We would like to thank our advisor, Elisabeti Kira, for her patience, trust, and
mentorship that enabled this project to be completed. We would also like to thank
the University of São Paulo for financing our research through the Programa
Unificado de Bolsas (PUB).

\vspace{0.5cm}
\bibliographystyle{plain}
\bibliography{references}

\end{document}
